% Options for packages loaded elsewhere
\PassOptionsToPackage{unicode}{hyperref}
\PassOptionsToPackage{hyphens}{url}
\PassOptionsToPackage{dvipsnames,svgnames,x11names}{xcolor}
\documentclass[
]{article}
\usepackage{xcolor}
\usepackage[margin = 0.5in]{geometry}
\usepackage{amsmath,amssymb}
\setcounter{secnumdepth}{-\maxdimen} % remove section numbering
\usepackage{iftex}
\ifPDFTeX
  \usepackage[T1]{fontenc}
  \usepackage[utf8]{inputenc}
  \usepackage{textcomp} % provide euro and other symbols
\else % if luatex or xetex
  \usepackage{unicode-math} % this also loads fontspec
  \defaultfontfeatures{Scale=MatchLowercase}
  \defaultfontfeatures[\rmfamily]{Ligatures=TeX,Scale=1}
\fi
\usepackage{lmodern}
\ifPDFTeX\else
  % xetex/luatex font selection
\fi
% Use upquote if available, for straight quotes in verbatim environments
\IfFileExists{upquote.sty}{\usepackage{upquote}}{}
\IfFileExists{microtype.sty}{% use microtype if available
  \usepackage[]{microtype}
  \UseMicrotypeSet[protrusion]{basicmath} % disable protrusion for tt fonts
}{}
\makeatletter
\@ifundefined{KOMAClassName}{% if non-KOMA class
  \IfFileExists{parskip.sty}{%
    \usepackage{parskip}
  }{% else
    \setlength{\parindent}{0pt}
    \setlength{\parskip}{6pt plus 2pt minus 1pt}}
}{% if KOMA class
  \KOMAoptions{parskip=half}}
\makeatother
\usepackage{graphicx}
\makeatletter
\newsavebox\pandoc@box
\newcommand*\pandocbounded[1]{% scales image to fit in text height/width
  \sbox\pandoc@box{#1}%
  \Gscale@div\@tempa{\textheight}{\dimexpr\ht\pandoc@box+\dp\pandoc@box\relax}%
  \Gscale@div\@tempb{\linewidth}{\wd\pandoc@box}%
  \ifdim\@tempb\p@<\@tempa\p@\let\@tempa\@tempb\fi% select the smaller of both
  \ifdim\@tempa\p@<\p@\scalebox{\@tempa}{\usebox\pandoc@box}%
  \else\usebox{\pandoc@box}%
  \fi%
}
% Set default figure placement to htbp
\def\fps@figure{htbp}
\makeatother
\setlength{\emergencystretch}{3em} % prevent overfull lines
\providecommand{\tightlist}{%
  \setlength{\itemsep}{0pt}\setlength{\parskip}{0pt}}
\usepackage{color}
\usepackage{framed}
\setlength{\fboxsep}{.8em}

\usepackage{tcolorbox}

\tcbset{
  mystyle/.style={
    before skip=-10pt, % Adjust this value
    after skip=-10pt   % Adjust this value
  }
}
\newtcolorbox{mynoteblock}[1][]{mystyle, #1}

\newtcolorbox{blackbox}{
  colback=black,
  colframe=orange,
  coltext=white,
  boxsep=5pt,
  arc=4pt}
  
\newtcolorbox{ltbluebox}{
  colback=blue!5!white,
  colframe=blue!75!black,
  boxsep=5pt,
  arc=4pt}
  
\newtcolorbox{ltgreenbox}{
  colback=green!5!white,
  colframe=blue!75!black,
  boxsep=5pt,
  arc=4pt}
  
\newtcolorbox{ltpurplebox}{
  colback=purple!5!white,
  colframe=blue!75!black,
  boxsep=5pt,
  arc=4pt}
  


\newenvironment{cols}[1][]{}{}

\newenvironment{col}[1]{\begin{minipage}[t]{#1}\ignorespaces\vspace{0pt}}{%
\end{minipage}
\ifhmode\unskip\fi
\aftergroup\useignorespacesandallpars}

\def\useignorespacesandallpars#1\ignorespaces\fi{%
#1\fi\ignorespacesandallpars}

\makeatletter
\def\ignorespacesandallpars{%
  \@ifnextchar\par
    {\expandafter\ignorespacesandallpars\@gobble}%
    {}%
}
\makeatother

\usepackage{awesomebox}
\setlength{\aweboxvskip}{1mm}


\pagenumbering{gobble}

\usepackage{enumitem}
\setlist{itemsep=1pt, parsep=0pt} %# Adjust vertical spacing
\setlist[itemize,1]{leftmargin=*, labelsep=0.3em} %# Adjust horizontal spacing
\setlist[enumerate,1]{leftmargin=*, labelsep=0.3em} %# Adjust horizontal spacing

\usepackage{bookmark}
\IfFileExists{xurl.sty}{\usepackage{xurl}}{} % add URL line breaks if available
\urlstyle{same}
\hypersetup{
  colorlinks=true,
  linkcolor={Maroon},
  filecolor={Maroon},
  citecolor={Blue},
  urlcolor={blue},
  pdfcreator={LaTeX via pandoc}}

\author{}
\date{\vspace{-2.5em}}

\begin{document}

\vspace{-1cm}

\section{Catch/Quota Setting Step 5: Quota
Setting}\label{catchquota-setting-step-5-quota-setting}

\begin{noteblock}
Catch and Quota Setting is a multistep process involving decisions by
the Northeast Regional Coordinating Committee (NRCC; including the Mid
Atlantic and New England Councils, Atlantic States Marine Fisheries
Commission, and NOAA's Greater Atlantic Regional Fisheries Office and
Northeast Fisheries Science Center), assessment scientists, assessment
peer reviewers, the Council's Scientific and Statistical Committee
(SSC), and the Council. Information for each stock and fishery is
reviewed and synthesized to determine catch and quota levels that meet
management objectives.

\end{noteblock}

\begin{itemize}
\tightlist
\item
  \textbf{Key question: What additional indicators could inform
  decisions at different steps in the Catch/Quota setting process? Could
  habitat, ecosystem, or socioeconomic indicators improve the process?}
\item
  \textbf{Scope: Should we develop a general process for using
  indicators in the catch/quota setting process, or focus on a specific
  application for a single species? Which species?}
\end{itemize}

\begin{cols}

\begin{col}{0.49\textwidth}

\begin{ltbluebox}

\subsection{What information is used
now?}\label{what-information-is-used-now}

\begin{ltpurplebox}
\textbf{5. Catch Limits, Targets, and Quotas}

The Council decides whether annually varying or averaged ABCs will be
used in multiyear harvest specifications, then allocates the ABC to
harvest sectors. Management uncertainty is determined and applied, then
discards are estimated and subtracted to derive quotas and harvest
limits.

\end{ltpurplebox}

\end{ltbluebox}

\end{col}

\begin{col}{0.02\textwidth}
~

\end{col}

\begin{col}{0.49\textwidth}

\begin{ltgreenbox}

\subsection{Information sources for
indicators}\label{information-sources-for-indicators}

\href{https://www.fisheries.noaa.gov/new-england-mid-atlantic/ecosystems/ecosystem-and-socioeconomic-profile-development-and-reports}{Ecosystem
and Socioeconomic Profiles} conceptual models and indicators for MAFMC
stocks:

\begin{itemize}
\tightlist
\item
  Black Sea Bass
\item
  Bluefish
\item
  Golden Tilefish
\item
  Longfin squid
\item
  Shortfin squid
\end{itemize}

\vspace{1ex}

\href{https://www.fisheries.noaa.gov/new-england-mid-atlantic/ecosystems/state-ecosystem-reports-northeast-us-shelf}{State
of the Ecosystem} and
\href{https://www.mafmc.org/s/2_MAB_RiskAssess_2024.pdf}{Risk
Assessment} indicators

\vspace{1ex}

The fish and invertebrate
\href{https://journals.plos.org/plosone/article?id=10.1371/journal.pone.0146756}{Climate
Vulnerability Analysis} sensitivity, exposure and vulnerability rankings

\vspace{1ex}

Multiple Essential Fish Habitat sources:

\begin{itemize}
\tightlist
\item
  \href{https://www.fisheries.noaa.gov/resource/map/essential-fish-habitat-mapper}{EFH
  Habitat Mapper}
\item
  \href{https://nrha.shinyapps.io/dataexplorer/}{NRHA Data Explorer}
\item
  \href{https://fishmaps.shinyapps.io/FishingEffectsDatabase/}{Fishing
  Gear Effects Database}
\end{itemize}

\vspace{1ex}

And many other potential sources:

\begin{itemize}
\tightlist
\item
  \href{https://www.fisheries.noaa.gov/inport/item/22549}{NEFSC Surveys}
\item
  \href{https://fwdp.shinyapps.io/tm2020/}{NEFSC Food Habits}
\item
  \href{https://www.vims.edu/research/units/programs/multispecies_fisheries_research/neamap/}{NEAMAP
  Survey and Food Habits}
\item
  \href{https://www.fisheries.noaa.gov/new-england-mid-atlantic/commercial-fishing/quota-monitoring-greater-atlantic-region}{GARFO
  Catch and Monitoring System}
\item
  \href{https://safis.accsp.org/accsp_prod/f?p=1490:1:6442641492940:::::}{Atlantic
  Coastal Cooperative Statistics Programa}
\item
  \href{https://www.fisheries.noaa.gov/recreational-fishing-data/introduction-marine-recreational-information-program-data}{Marine
  Recreational Information Program}
\item
  \href{https://data.marine.copernicus.eu/product/GLOBAL_MULTIYEAR_PHY_001_030/description}{GLORYS
  ocean reanalysis}
\end{itemize}

\subsubsection{Potential Pathways}\label{potential-pathways}

\begin{itemize}
\tightlist
\item
  Ecosystem information can enter at some portion of steps 1, 2, 4, and
  5. These pathways are detailed in separate 1 page summaries.
\item
  The Council could recommend additional assessment or review ToRs in
  steps 1 and 3.
\end{itemize}

\begin{ltpurplebox}

\subsection{Information gaps}\label{information-gaps}

Ecosystem and Socioeconomic Profiles evaluate and select stock specific
ecosystem and socioeconomic indicators that can be useful across many
steps of the Catch/Quota setting process; these are currently
unavailable for summer flounder, scup, mackerel, butterfish, surfclam,
ocean quahog, blueline tilefish, spiny dogfish, and monkfish.

\end{ltpurplebox}

\end{ltgreenbox}

\end{col}

\end{cols}

\end{document}
